\documentclass[aspectratio=169, 11pt]{beamer}
% ============================================================================
% Preambulo compartido para presentaciones Beamer
% Estadistica 2026-II - Pontificia Universidad Javeriana
% Incluir con: % ============================================================================
% Preambulo compartido para presentaciones Beamer
% Estadistica 2026-II - Pontificia Universidad Javeriana
% Incluir con: % ============================================================================
% Preambulo compartido para presentaciones Beamer
% Estadistica 2026-II - Pontificia Universidad Javeriana
% Incluir con: \input{../preambulo_beamer.tex} despues de \documentclass
% ============================================================================

\usepackage[T1]{fontenc}
\usepackage[utf8]{inputenc}
\usepackage[spanish]{babel}
\usepackage{amsmath, amssymb, amsfonts}
\usepackage{booktabs}
\usepackage{graphicx}
\usepackage{tikz}
\usetikzlibrary{shapes.geometric, positioning, arrows.meta, calc}
\usepackage{pgfplots}
\pgfplotsset{compat=1.18}
\usepackage{listings}
\usepackage{xcolor}
\usepackage{hyperref}
\usepackage{multicol}

% --- Tema Beamer (minimalista) ---
\usetheme{default}
\usecolortheme{default}
\usefonttheme{default}
\setbeamertemplate{navigation symbols}{}
\setbeamertemplate{caption}[numbered]
\setbeamertemplate{footline}{%
  \hspace{1em}{\tiny\url{https://github.com/polanco-jaime/Curso-Estadistica-2026-3}}\hfill\usebeamerfont{footline}\insertframenumber/\inserttotalframenumber\hspace{1em}\vspace{0.5em}%
}
\setbeamertemplate{itemize items}[circle]
\setbeamertemplate{enumerate items}[default]

% --- Colores institucionales (sutiles) ---
\definecolor{javeriana}{RGB}{0, 62, 105}
\definecolor{javerianalight}{RGB}{0, 105, 170}
\definecolor{codigobg}{RGB}{245, 245, 245}
\definecolor{codigofg}{RGB}{40, 40, 40}
\definecolor{comentario}{RGB}{100, 100, 100}
\definecolor{keyword}{RGB}{0, 100, 150}
\definecolor{string}{RGB}{180, 60, 30}

\setbeamercolor{title}{fg=javeriana}
\setbeamercolor{frametitle}{fg=javeriana}
\setbeamercolor{structure}{fg=javeriana}
\setbeamercolor{block title}{fg=white, bg=javeriana!75}
\setbeamercolor{block body}{bg=javeriana!5}
\setbeamercolor{block title alerted}{fg=white, bg=red!70!black}
\setbeamercolor{block body alerted}{bg=red!5}
\setbeamercolor{block title example}{fg=white, bg=green!40!black}
\setbeamercolor{block body example}{bg=green!5}
\setbeamercolor{item}{fg=javeriana}

\setbeamerfont{title}{series=\bfseries, size=\Large}
\setbeamerfont{frametitle}{series=\bfseries}

% --- Configuracion de listings para R ---
\lstset{
  language=R,
  basicstyle=\ttfamily\footnotesize,
  backgroundcolor=\color{codigobg},
  commentstyle=\color{comentario}\itshape,
  keywordstyle=\color{keyword}\bfseries,
  stringstyle=\color{string},
  numbers=none,
  frame=single,
  rulecolor=\color{javeriana!30},
  breaklines=true,
  showstringspaces=false,
  columns=flexible,
  morekeywords={library, ggplot, aes, geom_histogram, geom_boxplot,
    geom_point, geom_smooth, geom_density, geom_bar, geom_line,
    geom_vline, geom_hline, geom_errorbar, labs, theme_minimal,
    facet_wrap, scale_fill_manual, scale_color_manual,
    summarise, group_by, filter, mutate, select, arrange,
    read.csv, head, str, summary, mean, median, sd, var, cor,
    lm, t.test, chisq.test, wilcox.test, kruskal.test, aov,
    pnorm, qnorm, dnorm, rnorm, qt, pt, qchisq, pchisq, qf, pf,
    confint, coeftest, vcovHC, bptest, vif, cooks.distance,
    TukeyHSD, table, prop.table}
}

% --- Comandos personalizados ---
\newcommand{\R}{\texttt{R}}
\newcommand{\code}[1]{\texttt{\footnotesize #1}}
\newcommand{\variable}[1]{\texttt{#1}}
\newcommand{\paquete}[1]{\texttt{#1}}

% --- Informacion del curso ---
\institute{Pontificia Universidad Javeriana\\
  Facultad de Ciencias Econ\'omicas y Administrativas}
\date{2026-II}
 despues de \documentclass
% ============================================================================

\usepackage[T1]{fontenc}
\usepackage[utf8]{inputenc}
\usepackage[spanish]{babel}
\usepackage{amsmath, amssymb, amsfonts}
\usepackage{booktabs}
\usepackage{graphicx}
\usepackage{tikz}
\usetikzlibrary{shapes.geometric, positioning, arrows.meta, calc}
\usepackage{pgfplots}
\pgfplotsset{compat=1.18}
\usepackage{listings}
\usepackage{xcolor}
\usepackage{hyperref}
\usepackage{multicol}

% --- Tema Beamer (minimalista) ---
\usetheme{default}
\usecolortheme{default}
\usefonttheme{default}
\setbeamertemplate{navigation symbols}{}
\setbeamertemplate{caption}[numbered]
\setbeamertemplate{footline}{%
  \hspace{1em}{\tiny\url{https://github.com/polanco-jaime/Curso-Estadistica-2026-3}}\hfill\usebeamerfont{footline}\insertframenumber/\inserttotalframenumber\hspace{1em}\vspace{0.5em}%
}
\setbeamertemplate{itemize items}[circle]
\setbeamertemplate{enumerate items}[default]

% --- Colores institucionales (sutiles) ---
\definecolor{javeriana}{RGB}{0, 62, 105}
\definecolor{javerianalight}{RGB}{0, 105, 170}
\definecolor{codigobg}{RGB}{245, 245, 245}
\definecolor{codigofg}{RGB}{40, 40, 40}
\definecolor{comentario}{RGB}{100, 100, 100}
\definecolor{keyword}{RGB}{0, 100, 150}
\definecolor{string}{RGB}{180, 60, 30}

\setbeamercolor{title}{fg=javeriana}
\setbeamercolor{frametitle}{fg=javeriana}
\setbeamercolor{structure}{fg=javeriana}
\setbeamercolor{block title}{fg=white, bg=javeriana!75}
\setbeamercolor{block body}{bg=javeriana!5}
\setbeamercolor{block title alerted}{fg=white, bg=red!70!black}
\setbeamercolor{block body alerted}{bg=red!5}
\setbeamercolor{block title example}{fg=white, bg=green!40!black}
\setbeamercolor{block body example}{bg=green!5}
\setbeamercolor{item}{fg=javeriana}

\setbeamerfont{title}{series=\bfseries, size=\Large}
\setbeamerfont{frametitle}{series=\bfseries}

% --- Configuracion de listings para R ---
\lstset{
  language=R,
  basicstyle=\ttfamily\footnotesize,
  backgroundcolor=\color{codigobg},
  commentstyle=\color{comentario}\itshape,
  keywordstyle=\color{keyword}\bfseries,
  stringstyle=\color{string},
  numbers=none,
  frame=single,
  rulecolor=\color{javeriana!30},
  breaklines=true,
  showstringspaces=false,
  columns=flexible,
  morekeywords={library, ggplot, aes, geom_histogram, geom_boxplot,
    geom_point, geom_smooth, geom_density, geom_bar, geom_line,
    geom_vline, geom_hline, geom_errorbar, labs, theme_minimal,
    facet_wrap, scale_fill_manual, scale_color_manual,
    summarise, group_by, filter, mutate, select, arrange,
    read.csv, head, str, summary, mean, median, sd, var, cor,
    lm, t.test, chisq.test, wilcox.test, kruskal.test, aov,
    pnorm, qnorm, dnorm, rnorm, qt, pt, qchisq, pchisq, qf, pf,
    confint, coeftest, vcovHC, bptest, vif, cooks.distance,
    TukeyHSD, table, prop.table}
}

% --- Comandos personalizados ---
\newcommand{\R}{\texttt{R}}
\newcommand{\code}[1]{\texttt{\footnotesize #1}}
\newcommand{\variable}[1]{\texttt{#1}}
\newcommand{\paquete}[1]{\texttt{#1}}

% --- Informacion del curso ---
\institute{Pontificia Universidad Javeriana\\
  Facultad de Ciencias Econ\'omicas y Administrativas}
\date{2026-II}
 despues de \documentclass
% ============================================================================

\usepackage[T1]{fontenc}
\usepackage[utf8]{inputenc}
\usepackage[spanish]{babel}
\usepackage{amsmath, amssymb, amsfonts}
\usepackage{booktabs}
\usepackage{graphicx}
\usepackage{tikz}
\usetikzlibrary{shapes.geometric, positioning, arrows.meta, calc}
\usepackage{pgfplots}
\pgfplotsset{compat=1.18}
\usepackage{listings}
\usepackage{xcolor}
\usepackage{hyperref}
\usepackage{multicol}

% --- Tema Beamer (minimalista) ---
\usetheme{default}
\usecolortheme{default}
\usefonttheme{default}
\setbeamertemplate{navigation symbols}{}
\setbeamertemplate{caption}[numbered]
\setbeamertemplate{footline}{%
  \hspace{1em}{\tiny\url{https://github.com/polanco-jaime/Curso-Estadistica-2026-3}}\hfill\usebeamerfont{footline}\insertframenumber/\inserttotalframenumber\hspace{1em}\vspace{0.5em}%
}
\setbeamertemplate{itemize items}[circle]
\setbeamertemplate{enumerate items}[default]

% --- Colores institucionales (sutiles) ---
\definecolor{javeriana}{RGB}{0, 62, 105}
\definecolor{javerianalight}{RGB}{0, 105, 170}
\definecolor{codigobg}{RGB}{245, 245, 245}
\definecolor{codigofg}{RGB}{40, 40, 40}
\definecolor{comentario}{RGB}{100, 100, 100}
\definecolor{keyword}{RGB}{0, 100, 150}
\definecolor{string}{RGB}{180, 60, 30}

\setbeamercolor{title}{fg=javeriana}
\setbeamercolor{frametitle}{fg=javeriana}
\setbeamercolor{structure}{fg=javeriana}
\setbeamercolor{block title}{fg=white, bg=javeriana!75}
\setbeamercolor{block body}{bg=javeriana!5}
\setbeamercolor{block title alerted}{fg=white, bg=red!70!black}
\setbeamercolor{block body alerted}{bg=red!5}
\setbeamercolor{block title example}{fg=white, bg=green!40!black}
\setbeamercolor{block body example}{bg=green!5}
\setbeamercolor{item}{fg=javeriana}

\setbeamerfont{title}{series=\bfseries, size=\Large}
\setbeamerfont{frametitle}{series=\bfseries}

% --- Configuracion de listings para R ---
\lstset{
  language=R,
  basicstyle=\ttfamily\footnotesize,
  backgroundcolor=\color{codigobg},
  commentstyle=\color{comentario}\itshape,
  keywordstyle=\color{keyword}\bfseries,
  stringstyle=\color{string},
  numbers=none,
  frame=single,
  rulecolor=\color{javeriana!30},
  breaklines=true,
  showstringspaces=false,
  columns=flexible,
  morekeywords={library, ggplot, aes, geom_histogram, geom_boxplot,
    geom_point, geom_smooth, geom_density, geom_bar, geom_line,
    geom_vline, geom_hline, geom_errorbar, labs, theme_minimal,
    facet_wrap, scale_fill_manual, scale_color_manual,
    summarise, group_by, filter, mutate, select, arrange,
    read.csv, head, str, summary, mean, median, sd, var, cor,
    lm, t.test, chisq.test, wilcox.test, kruskal.test, aov,
    pnorm, qnorm, dnorm, rnorm, qt, pt, qchisq, pchisq, qf, pf,
    confint, coeftest, vcovHC, bptest, vif, cooks.distance,
    TukeyHSD, table, prop.table}
}

% --- Comandos personalizados ---
\newcommand{\R}{\texttt{R}}
\newcommand{\code}[1]{\texttt{\footnotesize #1}}
\newcommand{\variable}[1]{\texttt{#1}}
\newcommand{\paquete}[1]{\texttt{#1}}

% --- Informacion del curso ---
\institute{Pontificia Universidad Javeriana\\
  Facultad de Ciencias Econ\'omicas y Administrativas}
\date{2026-II}

\title{Sesión 4: Estimación Puntual e Intervalos de Confianza}
\subtitle{De estimadores puntuales a intervalos para medias y varianzas}
\author{Prof.\ Jaime Polanco Jim\'enez}
\date{Viernes 11 de septiembre de 2026}
\begin{document}

\begin{frame}
\titlepage
\end{frame}

\begin{frame}{Agenda}
\tableofcontents
\end{frame}

% ============================================================================
\section{Estimación Puntual}
% ============================================================================

\begin{frame}{Parámetros vs Estadísticos: Repaso}
\begin{block}{Conceptos fundamentales}
\begin{itemize}
  \item \textbf{Parámetro}: característica numérica de la \emph{población}
  \begin{itemize}
    \item Ejemplos: $\mu$ (media poblacional), $\sigma^2$ (varianza poblacional), $p$ (proporción poblacional)
    \item Usualmente \alert{desconocidos} (¿cuál es la media verdadera de inglés de TODOS los estudiantes de NI en Colombia?)
  \end{itemize}
  \item \textbf{Estadístico}: característica numérica de la \emph{muestra}
  \begin{itemize}
    \item Ejemplos: $\bar{x}$ (media muestral), $s^2$ (varianza muestral), $\hat{p}$ (proporción muestral)
    \item Son \alert{observables} (calculamos $\bar{x}$ con nuestra muestra)
  \end{itemize}
\end{itemize}
\end{block}

\vspace{0.3cm}
\textbf{Problema central de la inferencia}: usar estadísticos para aprender sobre parámetros
\end{frame}

\begin{frame}{¿Qué es la Estimación Puntual?}
\begin{definition}[Estimación puntual]
Es el proceso de usar un \textbf{estadístico muestral} para aproximar el valor de un \textbf{parámetro poblacional}.
\end{definition}

\vspace{0.3cm}
\begin{columns}[T]
\column{0.48\textwidth}
\textbf{Estimadores comunes}:
\begin{align*}
\bar{x} &\text{ estima } \mu \\
s^2 &\text{ estima } \sigma^2 \\
s &\text{ estima } \sigma \\
\hat{p} &\text{ estima } p
\end{align*}

\column{0.48\textwidth}
\textbf{Notación}:
\begin{itemize}
  \item $\hat{\theta}$ = estimador puntual del parámetro $\theta$
  \item Ejemplo: $\hat{\mu} = \bar{x}$
  \item $\hat{\sigma}^2 = s^2$
\end{itemize}
\end{columns}

\vspace{0.3cm}
\alert{Limitación}: un solo número no captura la incertidumbre de la estimación
\end{frame}

\begin{frame}{Propiedades Deseables de un Estimador}
\textbf{1. Insesgamiento (unbiasedness)}
\begin{itemize}
  \item Un estimador $\hat{\theta}$ es insesgado si $E(\hat{\theta}) = \theta$
  \item Su valor esperado (promedio a largo plazo) es igual al parámetro
  \item Ejemplos: $E(\bar{x}) = \mu$, $E(s^2) = \sigma^2$
\end{itemize}

\vspace{0.3cm}
\textbf{2. Eficiencia}
\begin{itemize}
  \item Entre dos estimadores insesgados, preferimos el de \alert{menor varianza}
  \item Menor varianza $\Rightarrow$ estimaciones más concentradas alrededor del parámetro
  \item Ejemplo: $\bar{x}$ es más eficiente que la mediana para estimar $\mu$ en distribuciones normales
\end{itemize}

\vspace{0.3cm}
\textbf{3. Consistencia}
\begin{itemize}
  \item A medida que $n \to \infty$, el estimador converge al parámetro
  \item $\bar{x} \xrightarrow{p} \mu$ cuando $n \to \infty$
\end{itemize}
\end{frame}

\begin{frame}[fragile]{Ejemplo: Estimación Puntual con Datos ICFES}
\textbf{Contexto}: Queremos estimar la media poblacional de \variable{MOD\_INGLES\_PUNT} para estudiantes de Negocios Internacionales en Colombia.

\vspace{0.3cm}
\begin{lstlisting}
# Cargar datos
icfes_ni <- read.csv("datos/icfes_ni_2023.csv")

# Estimacion puntual de mu
mu_estimado <- mean(icfes_ni$MOD_INGLES_PUNT, na.rm = TRUE)
cat("Estimacion puntual de mu:", round(mu_estimado, 2))

# Estimacion puntual de sigma^2
sigma2_estimado <- var(icfes_ni$MOD_INGLES_PUNT, na.rm = TRUE)
cat("\nEstimacion puntual de sigma^2:", round(sigma2_estimado, 2))

# Estimacion puntual de sigma
sigma_estimado <- sd(icfes_ni$MOD_INGLES_PUNT, na.rm = TRUE)
cat("\nEstimacion puntual de sigma:", round(sigma_estimado, 2))
\end{lstlisting}

\vspace{0.2cm}
\alert{Pregunta}: ¿Qué tan confiables son estas estimaciones? ¿Cuánto error hay?
\end{frame}

\begin{frame}{De Estimación Puntual a Intervalos}
\begin{columns}[c]
\column{0.48\textwidth}
\textbf{Estimación puntual}:
\begin{itemize}
  \item $\hat{\mu} = 162.5$ puntos
  \item Un solo valor
  \item No cuantifica incertidumbre
\end{itemize}

\vspace{0.5cm}
¿Es $\mu$ exactamente 162.5?\\
\alert{No}, pero probablemente está cerca.

\column{0.48\textwidth}
\textbf{Intervalo de confianza}:
\begin{itemize}
  \item $(159.8, 165.2)$ puntos
  \item Un rango de valores plausibles
  \item Cuantifica incertidumbre
\end{itemize}

\vspace{0.5cm}
Estamos \alert{95\% confiados} de que $\mu$ está en ese intervalo.
\end{columns}

\vspace{0.5cm}
\begin{block}{Transición}
Pasamos de un \textbf{punto} (estimación puntual) a un \textbf{intervalo} (estimación por intervalo) que refleja nuestra incertidumbre.
\end{block}
\end{frame}

% ============================================================================
\section{Lógica de los Intervalos de Confianza}
% ============================================================================

\begin{frame}{¿Qué es un Intervalo de Confianza?}
\begin{definition}[Intervalo de confianza]
Un \textbf{intervalo de confianza (IC)} para un parámetro $\theta$ al nivel $(1-\alpha) \times 100\%$ es un rango de valores que, si repitiéramos el muestreo muchas veces, contendría al parámetro verdadero en $(1-\alpha) \times 100\%$ de las muestras.
\end{definition}

\vspace{0.3cm}
\textbf{Estructura general}:
\[
\text{IC} = \text{Estimación puntual} \pm \text{Margen de error}
\]

\vspace{0.3cm}
\begin{block}{Nivel de confianza}
\begin{itemize}
  \item $(1-\alpha)$ = nivel de confianza (típicamente 0.90, 0.95, 0.99)
  \item $\alpha$ = nivel de significancia (típicamente 0.10, 0.05, 0.01)
  \item Ejemplo: IC al 95\% $\Rightarrow$ $(1-\alpha) = 0.95$, $\alpha = 0.05$
\end{itemize}
\end{block}
\end{frame}

\begin{frame}{Interpretación Correcta del IC}
\begin{alertblock}{Interpretación CORRECTA}
``Si tomáramos 100 muestras aleatorias diferentes y construyéramos un IC al 95\% con cada una, esperaríamos que aproximadamente \alert{95 de esos intervalos} contengan el parámetro verdadero $\mu$.''
\end{alertblock}

\vspace{0.3cm}
\begin{exampleblock}{Interpretación alternativa correcta}
``Estamos 95\% confiados de que el intervalo $(a, b)$ contiene a $\mu$.''
\end{exampleblock}

\vspace{0.3cm}
\begin{block}{Interpretación INCORRECTA}
\textcolor{red}{\textbf{NO}}: ``Hay un 95\% de probabilidad de que $\mu$ esté en $(a, b)$''.

\vspace{0.2cm}
\textbf{¿Por qué es incorrecta?} Porque $\mu$ es un valor \emph{fijo} (aunque desconocido), no una variable aleatoria. El intervalo es lo que varía de muestra a muestra.
\end{block}
\end{frame}

\begin{frame}{Visualización del Concepto de IC}
\begin{center}
\begin{tikzpicture}[scale=0.8]
  % Linea del parametro verdadero
  \draw[very thick, blue] (0,0) -- (10,0) node[right] {$\mu$ verdadero};
  \draw[blue, very thick] (5,0) circle (0.08);

  % Intervalos de confianza (simulados)
  \foreach \y/\xmin/\xmax/\color in {
    1/4.2/5.8/green!70!black,
    2/4.5/6.1/green!70!black,
    3/4.0/5.6/green!70!black,
    4/3.8/5.4/green!70!black,
    5/5.5/7.1/red,
    6/4.3/5.9/green!70!black,
    7/4.1/5.7/green!70!black,
    8/4.6/6.2/green!70!black}{
    \draw[\color, very thick] (\xmin,\y) -- (\xmax,\y);
    \draw[\color, fill=\color] (\xmin,\y) circle (0.04);
    \draw[\color, fill=\color] (\xmax,\y) circle (0.04);
    \pgfmathparse{(\xmin+\xmax)/2}
    \draw[\color, fill=white] (\pgfmathresult,\y) circle (0.06);
  }

  \node[anchor=west] at (10.2,4) {\small 8 muestras};
  \node[anchor=west, green!70!black] at (10.2,3) {\small 7 intervalos};
  \node[anchor=west, green!70!black] at (10.2,2.5) {\small contienen $\mu$};
  \node[anchor=west, red] at (10.2,2) {\small 1 intervalo};
  \node[anchor=west, red] at (10.2,1.5) {\small NO contiene $\mu$};
\end{tikzpicture}
\end{center}

\vspace{0.2cm}
\small
En este ejemplo, 7 de 8 IC (87.5\%) contienen a $\mu$. Con más muestras, converge a 95\%.
\end{frame}

\begin{frame}{Margen de Error}
\begin{definition}[Margen de error]
La mitad del ancho del intervalo de confianza. Para la media con $\sigma$ conocida:
\[
E = z_{\alpha/2} \cdot \frac{\sigma}{\sqrt{n}}
\]
\end{definition}

\vspace{0.3cm}
\textbf{Componentes}:
\begin{itemize}
  \item $z_{\alpha/2}$ = valor crítico de la distribución normal estándar
  \item $\frac{\sigma}{\sqrt{n}}$ = error estándar de $\bar{x}$
  \item Mayor confianza $(1-\alpha)$ $\Rightarrow$ mayor $z_{\alpha/2}$ $\Rightarrow$ mayor $E$
  \item Mayor $n$ $\Rightarrow$ menor $E$
\end{itemize}

\vspace{0.3cm}
\textbf{Trade-off}: Confianza vs Precisión
\begin{itemize}
  \item Más confianza (ej. 99\%) $\Rightarrow$ intervalo más ancho
  \item Menos confianza (ej. 90\%) $\Rightarrow$ intervalo más estrecho
\end{itemize}
\end{frame}

% ============================================================================
\section{IC para la Media: $\sigma$ Conocida}
% ============================================================================

\begin{frame}{IC para $\mu$ con $\sigma$ Conocida}
\textbf{Supuestos}:
\begin{itemize}
  \item Muestra aleatoria simple de tamaño $n$
  \item Población normal O $n \geq 30$ (TLC)
  \item Desviación estándar poblacional $\sigma$ \alert{conocida}
\end{itemize}

\vspace{0.3cm}
\begin{block}{Fórmula del IC al $(1-\alpha) \times 100\%$}
\[
\bar{x} \pm z_{\alpha/2} \frac{\sigma}{\sqrt{n}}
\]
Equivalentemente:
\[
\left( \bar{x} - z_{\alpha/2} \frac{\sigma}{\sqrt{n}}, \quad \bar{x} + z_{\alpha/2} \frac{\sigma}{\sqrt{n}} \right)
\]
\end{block}

\vspace{0.3cm}
\textbf{Fundamento}: Por el TLC, $\bar{x} \sim N\left(\mu, \frac{\sigma^2}{n}\right)$, entonces
\[
Z = \frac{\bar{x} - \mu}{\sigma/\sqrt{n}} \sim N(0,1)
\]
\end{frame}

\begin{frame}{Valores Críticos de $z$}
\begin{table}
\centering
\begin{tabular}{cccc}
\toprule
\textbf{Nivel de confianza} & $(1-\alpha)$ & $\alpha/2$ & $z_{\alpha/2}$ \\
\midrule
90\% & 0.90 & 0.05 & 1.645 \\
95\% & 0.95 & 0.025 & 1.960 \\
99\% & 0.99 & 0.005 & 2.576 \\
\bottomrule
\end{tabular}
\end{table}

\vspace{0.3cm}
\textbf{Interpretación}: Para un IC al 95\%, usamos $z_{0.025} = 1.96$, que deja 2.5\% en cada cola de la distribución normal estándar.

\vspace{0.3cm}
\begin{center}
\begin{tikzpicture}[scale=0.7]
  \draw[->] (-3.5,0) -- (3.5,0) node[right] {$z$};
  \draw[->] (0,0) -- (0,2.5) node[above] {};
  \draw[domain=-3.2:3.2, smooth, variable=\x, thick] plot ({\x},{2*exp(-0.5*\x*\x)});

  \fill[blue!20] (-1.96,0) -- plot[domain=-1.96:1.96] ({\x},{2*exp(-0.5*\x*\x)}) -- (1.96,0) -- cycle;
  \draw[dashed] (-1.96,0) -- (-1.96,{2*exp(-0.5*1.96*1.96)});
  \draw[dashed] (1.96,0) -- (1.96,{2*exp(-0.5*1.96*1.96)});

  \node[below] at (-1.96,0) {$-1.96$};
  \node[below] at (1.96,0) {$1.96$};
  \node at (0,1) {95\%};
  \node at (-2.7,0.3) {\small 2.5\%};
  \node at (2.7,0.3) {\small 2.5\%};
\end{tikzpicture}
\end{center}
\end{frame}

\begin{frame}[fragile]{Cálculo de Valores Críticos en R}
\begin{lstlisting}
# Valor critico para IC al 95%
z_95 <- qnorm(0.975)  # P(Z <= z) = 0.975
cat("z para 95%:", z_95)
# Resultado: 1.959964

# Valor critico para IC al 99%
z_99 <- qnorm(0.995)  # P(Z <= z) = 0.995
cat("\nz para 99%:", z_99)
# Resultado: 2.575829

# Valor critico para IC al 90%
z_90 <- qnorm(0.95)   # P(Z <= z) = 0.95
cat("\nz para 90%:", z_90)
# Resultado: 1.644854
\end{lstlisting}

\vspace{0.3cm}
\textbf{Nota}: Usamos \code{qnorm(1 - $\alpha$/2)} porque queremos el cuantil superior que deja $\alpha/2$ en la cola derecha.
\end{frame}

\begin{frame}[fragile]{Ejemplo Numérico: IC con $\sigma$ Conocida}
\textbf{Escenario}: Supongamos que de censos anteriores sabemos que $\sigma = 18$ para \variable{MOD\_INGLES\_PUNT} en NI. Tenemos una muestra de $n = 120$ estudiantes con $\bar{x} = 162.5$.

\vspace{0.3cm}
\textbf{Calcular IC al 95\%}:

\begin{lstlisting}
# Datos
xbar <- 162.5
sigma <- 18
n <- 120
confianza <- 0.95
alpha <- 1 - confianza

# Valor critico
z_critico <- qnorm(1 - alpha/2)

# Error estandar
ee <- sigma / sqrt(n)

# Margen de error
E <- z_critico * ee

# Limites del IC
ic_inferior <- xbar - E
ic_superior <- xbar + E

cat("IC al 95%: (", round(ic_inferior, 2), ",",
    round(ic_superior, 2), ")")
# IC al 95%: ( 159.28 , 165.72 )
\end{lstlisting}
\end{frame}

\begin{frame}{Interpretación del Resultado}
\begin{exampleblock}{Resultado}
IC al 95\% para $\mu$: $(159.28, 165.72)$ puntos
\end{exampleblock}

\vspace{0.3cm}
\textbf{Interpretación correcta}:
\begin{itemize}
  \item Estamos 95\% confiados de que la media poblacional de \variable{MOD\_INGLES\_PUNT} para estudiantes de NI está entre 159.28 y 165.72 puntos.
  \item Si repitiéramos el muestreo 100 veces, aproximadamente 95 de los intervalos construidos contendrían el verdadero valor de $\mu$.
\end{itemize}

\vspace{0.3cm}
\textbf{Implicación práctica}:
\begin{itemize}
  \item El margen de error es $E = 3.22$ puntos
  \item ¿Es suficientemente preciso para nuestros propósitos?
  \item Si necesitamos más precisión, podemos aumentar $n$
\end{itemize}
\end{frame}

% ============================================================================
\section{IC para la Media: $\sigma$ Desconocida (Distribución $t$)}
% ============================================================================

\begin{frame}{El Problema: $\sigma$ Desconocida}
\textbf{Situación real}: En la práctica, \alert{casi nunca} conocemos $\sigma$.

\vspace{0.3cm}
\textbf{Solución}: Usar la desviación estándar muestral $s$ en lugar de $\sigma$.

\vspace{0.3cm}
\begin{block}{Consecuencia}
Al reemplazar $\sigma$ por $s$, introducimos una fuente adicional de variabilidad. El estadístico
\[
t = \frac{\bar{x} - \mu}{s/\sqrt{n}}
\]
ya NO sigue una distribución normal estándar, sino una \alert{distribución $t$ de Student} con $n-1$ grados de libertad.
\end{block}

\vspace{0.3cm}
\textbf{Notación}: $t \sim t(n-1)$ o $t_{n-1}$
\end{frame}

\begin{frame}{Distribución $t$ de Student}
\begin{columns}[T]
\column{0.48\textwidth}
\textbf{Características}:
\begin{itemize}
  \item Simétrica, centrada en 0
  \item Forma de campana (similar a $N(0,1)$)
  \item \alert{Colas más pesadas} que la normal
  \item Depende de los grados de libertad (gl = $n-1$)
  \item Con $n$ grande, $t(n-1) \approx N(0,1)$
\end{itemize}

\column{0.48\textwidth}
\begin{tikzpicture}[scale=0.65]
  \draw[->] (-3.5,0) -- (3.5,0) node[right] {$t$};
  \draw[->] (0,0) -- (0,2.5);

  % Normal estandar
  \draw[domain=-3.2:3.2, smooth, variable=\x, thick, blue]
    plot ({\x},{2*exp(-0.5*\x*\x)}) node[above right] {\small $N(0,1)$};

  % t con gl=5
  \draw[domain=-3.2:3.2, smooth, variable=\x, thick, red, dashed]
    plot ({\x},{2*0.92*(1+\x*\x/5)^(-3)}) node[above left] {\small $t(5)$};

  \node[below] at (0,0) {0};
\end{tikzpicture}

\vspace{0.2cm}
\small La $t$ con pocos gl tiene colas más anchas.
\end{columns}

\vspace{0.3cm}
\textbf{Implicación}: Para el mismo nivel de confianza, los valores críticos $t_{\alpha/2, n-1}$ son \alert{mayores} que $z_{\alpha/2}$, lo que resulta en intervalos \alert{más anchos}.
\end{frame}

\begin{frame}{IC para $\mu$ con $\sigma$ Desconocida}
\textbf{Supuestos}:
\begin{itemize}
  \item Muestra aleatoria simple de tamaño $n$
  \item Población normal (más importante cuando $n$ es pequeño)
  \item $\sigma$ \alert{desconocida}
\end{itemize}

\vspace{0.3cm}
\begin{block}{Fórmula del IC al $(1-\alpha) \times 100\%$}
\[
\bar{x} \pm t_{\alpha/2, n-1} \frac{s}{\sqrt{n}}
\]
donde:
\begin{itemize}
  \item $t_{\alpha/2, n-1}$ = valor crítico de la distribución $t$ con $n-1$ grados de libertad
  \item $s$ = desviación estándar muestral
  \item $\frac{s}{\sqrt{n}}$ = error estándar estimado de $\bar{x}$
\end{itemize}
\end{block}
\end{frame}

\begin{frame}{Comparación de Valores Críticos: $z$ vs $t$}
\begin{table}
\centering
\small
\begin{tabular}{ccccc}
\toprule
\textbf{gl} & $n$ & $t_{0.025, n-1}$ & $z_{0.025}$ & Diferencia \\
\midrule
5  & 6   & 2.571 & 1.960 & +31\% \\
10 & 11  & 2.228 & 1.960 & +14\% \\
20 & 21  & 2.086 & 1.960 & +6\% \\
30 & 31  & 2.042 & 1.960 & +4\% \\
60 & 61  & 2.000 & 1.960 & +2\% \\
$\infty$ & $\infty$ & 1.960 & 1.960 & 0\% \\
\bottomrule
\end{tabular}
\end{table}

\vspace{0.3cm}
\textbf{Observación}: Con $n \geq 30$, la diferencia es pequeña ($<5\%$), pero sigue siendo mejor usar $t$ cuando $\sigma$ es desconocida.
\end{frame}

\begin{frame}[fragile]{Cálculo de Valores Críticos $t$ en R}
\begin{lstlisting}
# Valor critico t para IC al 95% con diferentes gl
gl_valores <- c(5, 10, 20, 30, 60, 120)
alpha <- 0.05

resultados <- data.frame(
  gl = gl_valores,
  n = gl_valores + 1,
  t_critico = qt(1 - alpha/2, df = gl_valores),
  z_critico = qnorm(1 - alpha/2)
)

resultados$diferencia <-
  (resultados$t_critico / resultados$z_critico - 1) * 100

print(round(resultados, 3))
\end{lstlisting}

\vspace{0.3cm}
\textbf{Función clave}: \code{qt(p, df)} devuelve el cuantil $p$ de la distribución $t$ con \code{df} grados de libertad.
\end{frame}

\begin{frame}[fragile]{Ejemplo: IC para la Media de MOD\_INGLES\_PUNT en NI}
\textbf{Datos reales}: Muestra de $n = 120$ estudiantes de NI.

\begin{lstlisting}
# Cargar datos
icfes_ni <- read.csv("datos/icfes_ni_2023.csv")

# Calcular estadisticos muestrales
ingles <- icfes_ni$MOD_INGLES_PUNT
n <- length(ingles)
xbar <- mean(ingles, na.rm = TRUE)
s <- sd(ingles, na.rm = TRUE)

# IC al 95% usando la formula
alpha <- 0.05
gl <- n - 1
t_critico <- qt(1 - alpha/2, df = gl)
ee <- s / sqrt(n)
E <- t_critico * ee

ic_inf <- xbar - E
ic_sup <- xbar + E

cat("n =", n, "\nxbar =", round(xbar, 2),
    "\ns =", round(s, 2))
cat("\nIC al 95%: (", round(ic_inf, 2), ",",
    round(ic_sup, 2), ")")
\end{lstlisting}
\end{frame}

\begin{frame}[fragile]{Método Directo con \code{t.test()}}
\R\ tiene una función integrada que calcula el IC automáticamente:

\begin{lstlisting}
# Metodo directo: t.test()
resultado <- t.test(ingles, conf.level = 0.95)

# Ver el resultado completo
print(resultado)

# Extraer solo el IC
ic <- resultado$conf.int
cat("IC al 95%:", round(ic, 2))

# Cambiar el nivel de confianza
ic_99 <- t.test(ingles, conf.level = 0.99)$conf.int
cat("\nIC al 99%:", round(ic_99, 2))

ic_90 <- t.test(ingles, conf.level = 0.90)$conf.int
cat("\nIC al 90%:", round(ic_90, 2))
\end{lstlisting}

\vspace{0.3cm}
\textbf{Ventaja}: \code{t.test()} también realiza una prueba de hipótesis (que veremos en sesiones futuras).
\end{frame}

\begin{frame}[fragile]{Interpretación del Output de \code{t.test()}}
\begin{lstlisting}[basicstyle=\ttfamily\footnotesize]
	One Sample t-test

data:  ingles
t = 42.356, df = 119, p-value < 2.2e-16
alternative hypothesis: true mean is not equal to 0
95 percent confidence interval:
 159.28 165.72
sample estimates:
mean of x
   162.50
\end{lstlisting}

\vspace{0.3cm}
\textbf{Elementos clave}:
\begin{itemize}
  \item \textbf{t = 42.356}: estadístico $t$ para prueba $H_0: \mu = 0$ (no relevante aquí)
  \item \textbf{df = 119}: grados de libertad ($n-1$)
  \item \textbf{95\% CI: (159.28, 165.72)}: ¡el intervalo de confianza!
  \item \textbf{mean of x = 162.50}: estimación puntual $\bar{x}$
\end{itemize}
\end{frame}

% ============================================================================
\section{IC para la Varianza (Distribución $\chi^2$)}
% ============================================================================

\begin{frame}{IC para la Varianza $\sigma^2$}
\textbf{Motivación}: A veces queremos estimar la \alert{variabilidad} poblacional, no solo la media.

\vspace{0.3cm}
\textbf{Aplicaciones}:
\begin{itemize}
  \item Educación: ¿qué tan heterogéneo es el rendimiento en inglés entre estudiantes de NI?
  \item Comercio internacional: volatilidad de la TRM (tasa de cambio) para planear importaciones
  \item Logística: variabilidad en tiempos de entrega de MercadoLibre o Rappi
\end{itemize}

\vspace{0.3cm}
\begin{block}{Distribución muestral de $s^2$}
Si la población es normal, entonces:
\[
\frac{(n-1)s^2}{\sigma^2} \sim \chi^2(n-1)
\]
donde $\chi^2(n-1)$ es la distribución chi-cuadrado con $n-1$ grados de libertad.
\end{block}
\end{frame}

\begin{frame}{Distribución Chi-Cuadrado ($\chi^2$)}
\begin{columns}[T]
\column{0.48\textwidth}
\textbf{Características}:
\begin{itemize}
  \item \alert{Asimétrica} (sesgada a la derecha)
  \item Solo valores positivos ($\chi^2 \geq 0$)
  \item Depende de los gl
  \item Con gl grandes, se aproxima a la normal
\end{itemize}

\vspace{0.3cm}
\textbf{Implicación}: El IC para $\sigma^2$ será \alert{asimétrico}.

\column{0.48\textwidth}
\begin{tikzpicture}[scale=0.65]
  \draw[->] (-0.3,0) -- (10,0) node[right] {$\chi^2$};
  \draw[->] (0,0) -- (0,3);

  % Chi-cuadrado con gl=3
  \draw[domain=0.1:9, smooth, variable=\x, thick, blue, samples=100]
    plot ({\x},{2.5*(\x^0.5)*exp(-\x/2)}) node[right] {\small $\chi^2(3)$};

  % Chi-cuadrado con gl=6
  \draw[domain=0.1:9, smooth, variable=\x, thick, red, dashed, samples=100]
    plot ({\x},{1.8*(\x^1)*exp(-\x/2)}) node[above right] {\small $\chi^2(6)$};

  \node[below] at (0,0) {0};
\end{tikzpicture}
\end{columns}

\vspace{0.3cm}
\small
\textbf{Nota}: A diferencia de $z$ y $t$, la $\chi^2$ NO es simétrica, así que los dos valores críticos no son opuestos.
\end{frame}

\begin{frame}{Fórmula del IC para $\sigma^2$}
\textbf{Supuestos}:
\begin{itemize}
  \item Muestra aleatoria simple de tamaño $n$
  \item Población \alert{normal} (más importante que para la media)
\end{itemize}

\vspace{0.3cm}
\begin{block}{IC al $(1-\alpha) \times 100\%$ para $\sigma^2$}
\[
\left( \frac{(n-1)s^2}{\chi^2_{\alpha/2, n-1}}, \quad \frac{(n-1)s^2}{\chi^2_{1-\alpha/2, n-1}} \right)
\]
\end{block}

\vspace{0.3cm}
\textbf{Valores críticos}:
\begin{itemize}
  \item $\chi^2_{\alpha/2, n-1}$ = valor que deja $\alpha/2$ en la cola \alert{derecha} (cuantil superior)
  \item $\chi^2_{1-\alpha/2, n-1}$ = valor que deja $\alpha/2$ en la cola \alert{izquierda} (cuantil inferior)
\end{itemize}

\vspace{0.3cm}
\textbf{IC para $\sigma$}: Tomar raíz cuadrada de los límites.
\end{frame}

\begin{frame}[fragile]{Ejemplo: IC para la Varianza de MOD\_INGLES\_PUNT}
\begin{lstlisting}
# Datos
n <- 120
s2 <- var(ingles, na.rm = TRUE)
gl <- n - 1
alpha <- 0.05

# Valores criticos chi-cuadrado
chi2_inf <- qchisq(alpha/2, df = gl)       # Cola izquierda
chi2_sup <- qchisq(1 - alpha/2, df = gl)   # Cola derecha

# IC para sigma^2
ic_var_inf <- (gl * s2) / chi2_sup
ic_var_sup <- (gl * s2) / chi2_inf

cat("s^2 =", round(s2, 2))
cat("\nIC al 95% para sigma^2: (",
    round(ic_var_inf, 2), ",", round(ic_var_sup, 2), ")")

# IC para sigma (tomar raiz cuadrada)
ic_sd_inf <- sqrt(ic_var_inf)
ic_sd_sup <- sqrt(ic_var_sup)

cat("\nIC al 95% para sigma: (",
    round(ic_sd_inf, 2), ",", round(ic_sd_sup, 2), ")")
\end{lstlisting}
\end{frame}

\begin{frame}{Asimetría del IC para $\sigma^2$}
\textbf{Ejemplo numérico}: Supongamos $n = 120$, $s^2 = 324$ (entonces $s = 18$).

\vspace{0.3cm}
\begin{itemize}
  \item $\chi^2_{0.025, 119} = 152.22$ (cola derecha)
  \item $\chi^2_{0.975, 119} = 89.39$ (cola izquierda)
\end{itemize}

\vspace{0.3cm}
IC para $\sigma^2$:
\[
\left( \frac{119 \times 324}{152.22}, \quad \frac{119 \times 324}{89.39} \right) = (253.5, 431.6)
\]

\vspace{0.3cm}
IC para $\sigma$:
\[
(\sqrt{253.5}, \sqrt{431.6}) = (15.9, 20.8)
\]

\vspace{0.3cm}
\textbf{Observación}: La estimación puntual $s = 18$ NO está en el centro del intervalo (15.9, 20.8) debido a la asimetría de la distribución $\chi^2$.
\end{frame}

% ============================================================================
\section{Forest Plot Departamental}
% ============================================================================

\begin{frame}{IC por Departamento: Motivación}
\textbf{Pregunta}: ¿El nivel de inglés varía significativamente entre departamentos de Colombia?

\vspace{0.3cm}
\textbf{Estrategia}:
\begin{enumerate}
  \item Calcular IC al 95\% para $\mu$ en cada departamento
  \item Visualizar todos los IC en un \alert{forest plot}
  \item Comparar: si dos intervalos NO se solapan, hay evidencia de diferencia
\end{enumerate}

\vspace{0.3cm}
\textbf{Aplicación}: Identificar regiones con rendimiento excepcionalmente alto o bajo en inglés para diseñar políticas educativas focalizadas.
\end{frame}

\begin{frame}[fragile]{Cálculo de IC por Departamento en R}
\begin{lstlisting}
library(dplyr)

# Calcular IC por departamento
ic_deptos <- icfes_ni %>%
  group_by(ESTU_DEPTO_RESIDE) %>%
  summarise(
    n = n(),
    media = mean(MOD_INGLES_PUNT, na.rm = TRUE),
    sd = sd(MOD_INGLES_PUNT, na.rm = TRUE),
    ee = sd / sqrt(n),
    gl = n - 1,
    t_critico = qt(0.975, df = gl),
    ic_inf = media - t_critico * ee,
    ic_sup = media + t_critico * ee
  ) %>%
  filter(n >= 30)  # Solo deptos con muestra suficiente

# Ver primeros resultados
head(ic_deptos)
\end{lstlisting}
\end{frame}

\begin{frame}[fragile]{Crear Forest Plot con ggplot2}
\begin{lstlisting}
library(ggplot2)

# Ordenar por media para mejor visualizacion
ic_deptos <- ic_deptos %>%
  arrange(media)

# Forest plot
ggplot(ic_deptos, aes(x = media, y = reorder(ESTU_DEPTO_RESIDE, media))) +
  geom_point(size = 3, color = "steelblue") +
  geom_errorbarh(aes(xmin = ic_inf, xmax = ic_sup),
                 height = 0.3, color = "steelblue") +
  geom_vline(xintercept = mean(ic_deptos$media),
             linetype = "dashed", color = "red") +
  labs(x = "Media de MOD_INGLES_PUNT (IC 95%)",
       y = "Departamento",
       title = "Intervalos de Confianza por Departamento",
       subtitle = "Estudiantes de Negocios Internacionales") +
  theme_minimal()
\end{lstlisting}
\end{frame}

\begin{frame}{Interpretación del Forest Plot}
\textbf{Elementos visuales}:
\begin{itemize}
  \item \textbf{Punto}: estimación puntual $\bar{x}$ para cada departamento
  \item \textbf{Barra horizontal}: IC al 95\%
  \item \textbf{Línea vertical roja}: media nacional (de los departamentos con $n \geq 30$)
\end{itemize}

\vspace{0.3cm}
\textbf{Interpretación}:
\begin{itemize}
  \item Departamentos con IC que NO cruzan la línea roja son significativamente diferentes de la media nacional
  \item Si dos IC NO se solapan, hay evidencia de diferencia entre esos departamentos
  \item IC más anchos $\Rightarrow$ muestras más pequeñas o mayor variabilidad
\end{itemize}

\vspace{0.3cm}
\textbf{Ejemplo hipotético}:
\begin{itemize}
  \item Bogotá: IC = (168, 172) $\Rightarrow$ significativamente superior
  \item Chocó: IC = (140, 148) $\Rightarrow$ significativamente inferior
\end{itemize}
\end{frame}

\begin{frame}[fragile]{Identificar Departamentos Atípicos}
\begin{lstlisting}
# Media nacional (ponderada por n)
media_nacional <- weighted.mean(ic_deptos$media, ic_deptos$n)

# Departamentos significativamente superiores
superiores <- ic_deptos %>%
  filter(ic_inf > media_nacional) %>%
  select(ESTU_DEPTO_RESIDE, media, ic_inf, ic_sup)

cat("Departamentos significativamente superiores:\n")
print(superiores)

# Departamentos significativamente inferiores
inferiores <- ic_deptos %>%
  filter(ic_sup < media_nacional) %>%
  select(ESTU_DEPTO_RESIDE, media, ic_inf, ic_sup)

cat("\nDepartamentos significativamente inferiores:\n")
print(inferiores)
\end{lstlisting}
\end{frame}

\begin{frame}[fragile]{Comparación Visual: Ancho de los IC}
\begin{lstlisting}
# Calcular ancho del IC
ic_deptos <- ic_deptos %>%
  mutate(ancho_ic = ic_sup - ic_inf)

# Grafico: ancho vs tamano de muestra
ggplot(ic_deptos, aes(x = n, y = ancho_ic)) +
  geom_point(size = 3, alpha = 0.6) +
  geom_smooth(method = "lm", se = FALSE,
              color = "red", linetype = "dashed") +
  labs(x = "Tamano de muestra (n)",
       y = "Ancho del IC",
       title = "Relacion entre n y Precision",
       subtitle = "Ancho del IC = f(1/sqrt(n))") +
  theme_minimal()
\end{lstlisting}

\vspace{0.3cm}
\textbf{Relación esperada}: $\text{Ancho} = 2 \times t_{\alpha/2} \times \frac{s}{\sqrt{n}}$, entonces el ancho $\propto 1/\sqrt{n}$.
\end{frame}

% ============================================================================
\section{Resumen: ¿Cuándo usar $z$, $t$, $\chi^2$?}
% ============================================================================

\begin{frame}{Tabla Resumen de Intervalos de Confianza}
\begin{table}
\centering
\footnotesize
\begin{tabular}{p{2.5cm}p{2.5cm}p{3cm}p{3cm}}
\toprule
\textbf{Parámetro} & \textbf{Condición} & \textbf{Distribución} & \textbf{Fórmula} \\
\midrule
$\mu$ & $\sigma$ conocida & $z$ (normal) & $\bar{x} \pm z_{\alpha/2} \frac{\sigma}{\sqrt{n}}$ \\[0.3cm]
$\mu$ & $\sigma$ desconocida & $t(n-1)$ & $\bar{x} \pm t_{\alpha/2,n-1} \frac{s}{\sqrt{n}}$ \\[0.3cm]
$\sigma^2$ & Población normal & $\chi^2(n-1)$ & $\left(\frac{(n-1)s^2}{\chi^2_{\alpha/2}}, \frac{(n-1)s^2}{\chi^2_{1-\alpha/2}}\right)$ \\
\bottomrule
\end{tabular}
\end{table}

\vspace{0.3cm}
\textbf{Nota}: En la práctica, casi siempre usamos la distribución $t$ para la media, porque $\sigma$ es desconocida.
\end{frame}

\begin{frame}{Diagrama de Flujo para Elegir el IC}
\begin{center}
\begin{tikzpicture}[
  node distance=1.2cm and 1.5cm,
  decision/.style={diamond, draw, fill=blue!20, text width=3cm, align=center, inner sep=0pt},
  block/.style={rectangle, draw, fill=green!20, text width=3.5cm, align=center, rounded corners, minimum height=1cm},
  line/.style={draw, -{Latex[length=2mm]}}
]

  \node [block] (inicio) {?`Qu\'e par\'ametro estimar?};

  \node [decision, below=0.8cm of inicio] (param) {$\mu$ o $\sigma^2$?};

  \node [decision, below left=1.2cm and 2cm of param] (sigma_conocida) {?`$\sigma$ conocida?};
  \node [block, below=1cm of sigma_conocida] (usa_z) {Usar distribuci\'on $z$};
  \node [block, right=2cm of usa_z] (usa_t) {Usar distribuci\'on $t$};

  \node [block, below right=1.2cm and 2cm of param] (usa_chi2) {Usar distribuci\'on $\chi^2$};

  \path [line] (inicio) -- (param);
  \path [line] (param) -- node[above, sloped] {\small $\mu$} (sigma_conocida);
  \path [line] (param) -- node[above, sloped] {\small $\sigma^2$} (usa_chi2);
  \path [line] (sigma_conocida) -- node[left] {\small S\'i} (usa_z);
  \path [line] (sigma_conocida) -- node[above] {\small No} (usa_t);

\end{tikzpicture}
\end{center}

\vspace{0.3cm}
\small
\textbf{Regla práctica}: Si estás estimando $\mu$ y no te dicen explícitamente que $\sigma$ es conocida, \alert{usa la distribución $t$}.
\end{frame}

\begin{frame}{Supuestos y Robustez}
\begin{table}
\centering
\small
\begin{tabular}{lp{5cm}p{5cm}}
\toprule
\textbf{Método} & \textbf{Supuesto clave} & \textbf{Robustez} \\
\midrule
IC para $\mu$ (z o t) & Normalidad de la población O $n \geq 30$ (TLC) & \textbf{Robusta} con $n$ grande incluso si la población no es normal \\[0.3cm]
IC para $\sigma^2$ ($\chi^2$) & Población \alert{normal} & \textbf{NO robusta}. Sensible a desviaciones de normalidad \\
\bottomrule
\end{tabular}
\end{table}

\vspace{0.3cm}
\textbf{Recomendaciones}:
\begin{itemize}
  \item Para $\mu$: si $n < 30$, verificar normalidad con QQ-plot o Shapiro-Wilk
  \item Para $\sigma^2$: siempre verificar normalidad, considerar métodos bootstrap si hay dudas
\end{itemize}
\end{frame}

% ============================================================================
\section{Conclusiones}
% ============================================================================

\begin{frame}{Resumen de la Sesión}
\textbf{Hoy aprendimos}:
\begin{enumerate}
  \item \textbf{Estimación puntual}: usar $\bar{x}$ para estimar $\mu$, $s^2$ para estimar $\sigma^2$
  \item \textbf{Intervalos de confianza}: cuantificar incertidumbre alrededor de la estimación
  \item \textbf{IC para $\mu$}:
  \begin{itemize}
    \item Con $\sigma$ conocida: distribución $z$
    \item Con $\sigma$ desconocida: distribución $t$ (caso real)
  \end{itemize}
  \item \textbf{IC para $\sigma^2$}: distribución $\chi^2$ (asimétrico)
  \item \textbf{Forest plots}: comparar IC entre grupos (departamentos)
\end{enumerate}

\vspace{0.3cm}
\textbf{Próxima sesión}: IC para proporciones, tamaño de muestra, y muestreo estratificado.
\end{frame}

\begin{frame}{Ejercicios para la Próxima Clase}
\textbf{Con los datos de ICFES}:
\begin{enumerate}
  \item Calcular IC al 95\% y 99\% para la media de \variable{MOD\_RAZONA\_CUANTITAT} en NI
  \item Calcular IC al 95\% para la desviación estándar de \variable{MOD\_INGLES\_PUNT}
  \item Crear un forest plot comparando IC de inglés entre tipo de IES (pública vs privada)
  \item Investigar: ¿Cuántos departamentos tienen IC que NO se solapan con Bogotá?
  \item \textbf{Desafío}: Calcular el IC para la \emph{mediana} usando bootstrap (investigar método)
\end{enumerate}

\vspace{0.3cm}
\textbf{Lectura}: Anderson et al., Capítulo 8 (Intervalos de Confianza)
\end{frame}

\begin{frame}
\begin{center}
\Large
¡Gracias por su atención!

\vspace{1cm}
\normalsize
Preguntas y discusión
\end{center}
\end{frame}

\end{document}
